% !TEX program = xelatex
% Roteiro de Estudos — Realidade de Empresas Produtoras de Soluções de IA
% Tema escuro (preto e verde), duas colunas — compilar com XeLaTeX
\documentclass[10pt]{article}

% ---------- Fontes / codificação ----------
\usepackage{fontspec}
\usepackage{amsmath}
% Fontes padrão do LaTeX (família Latin Modern) — já vêm em qualquer
% distribuição (MiKTeX/TeX Live), não precisam ser instaladas separadamente.
% Usamos a versão SEM SERIFA (Latin Modern Sans), mais confortável na tela.
% Para hifenização PT-BR ao compilar localmente, descomente a linha a seguir
% (o MiKTeX/TeX Live baixam o pacote de português automaticamente):
% \usepackage[brazilian]{babel}
\setmainfont{Latin Modern Sans}
\newfontfamily\displayfont{Latin Modern Sans}
\setmonofont{Latin Modern Mono}

% ---------- Layout ----------
\usepackage[a4paper,top=1.55cm,bottom=1.75cm,left=1.45cm,right=1.45cm]{geometry}
\usepackage{multicol}
\usepackage{microtype}
\usepackage{setspace}
\usepackage{etoolbox}
\onehalfspacing % espaçamento 1,5 entre linhas nos textos
% Mantém as TABELAS com espaçamento simples (senão as células ficam altas demais)
\AtBeginEnvironment{tabularx}{\singlespacing}
\usepackage{ragged2e}
\usepackage{needspace}

% ---------- Cores ----------
\usepackage{xcolor}
\definecolor{bg}{HTML}{0B0F0D}        % fundo quase preto
\definecolor{panel}{HTML}{141B18}     % cartões / linhas alternadas
\definecolor{panelHi}{HTML}{1B2521}   % painel mais claro
\definecolor{accent}{HTML}{33E086}    % verde vivo (destaque)
\definecolor{accentDeep}{HTML}{0F5A3A}% verde profundo (cabeçalho de tabela)
\definecolor{accentDim}{HTML}{1E8F5E} % verde médio
\definecolor{textcol}{HTML}{E8F0EC}   % texto principal
\definecolor{muted}{HTML}{93A79D}     % legendas / fontes / apoio
\definecolor{line}{HTML}{27362F}      % separadores sutis

\usepackage{pagecolor}
\pagecolor{bg}

% ---------- Tabelas ----------
\usepackage{array}
\usepackage{tabularx}
\usepackage{booktabs}
\usepackage{colortbl}
\arrayrulecolor{line}
\renewcommand{\arraystretch}{1.28}
\newcolumntype{L}{>{\RaggedRight\arraybackslash}X}
\newcolumntype{Y}{>{\RaggedRight\arraybackslash\columncolor{panel}}X}
% Colunas ponderadas para tabelas-grade de responsabilidades
\newcolumntype{R}{>{\hsize=1.34\hsize\RaggedRight\arraybackslash}X}
\newcolumntype{I}{>{\hsize=0.66\hsize\RaggedRight\arraybackslash}X}
% Lista compacta para dentro de células de tabela
\newenvironment{cbul}
  {\begin{itemize}[leftmargin=1.05em,topsep=1pt,itemsep=1.5pt,parsep=0pt,after=\vspace{1pt}]}
  {\end{itemize}}

% ---------- Caixas ----------
\usepackage{tcolorbox}
\tcbuselibrary{skins}

% ---------- Listas ----------
\usepackage{enumitem}
\newcommand{\sqbul}{\textcolor{accent}{\raisebox{0.12ex}{\rule{0.42em}{0.42em}}}}
\setlist[itemize,1]{leftmargin=1.35em,itemsep=2.5pt,topsep=3pt,parsep=0pt,
  label=\sqbul}
\setlist[itemize,2]{leftmargin=1.2em,itemsep=1.5pt,topsep=2pt,
  label=\textcolor{accentDim}{\textbullet}}
\setlist[enumerate,1]{leftmargin=1.7em,itemsep=2.5pt,topsep=3pt,parsep=0pt,
  label=\textcolor{accent}{\bfseries\arabic*.}}

% ---------- Imagens ----------
\usepackage{graphicx}
\usepackage{adjustbox}
\graphicspath{{images/}}

% ---------- Sumário (tocloft) ----------
\usepackage{tocloft}
\renewcommand{\contentsname}{Conteúdo}
\renewcommand{\cfttoctitlefont}{\displayfont\Huge\bfseries\color{accent}}
\renewcommand{\cftsecfont}{\displayfont\bfseries\color{accent}}
\renewcommand{\cftsecpagefont}{\bfseries\color{accent}}
\renewcommand{\cftsubsecfont}{\color{textcol}}
\renewcommand{\cftsubsecpagefont}{\color{muted}}
\renewcommand{\cftsubsubsecfont}{\color{muted}}
\renewcommand{\cftsubsubsecpagefont}{\color{muted}}
\renewcommand{\cftsecleader}{\hfill}
\renewcommand{\cftsubsecleader}{\color{line}\cftdotfill{\cftsubsecdotsep}}
\renewcommand{\cftsubsubsecleader}{\color{line}\cftdotfill{\cftsubsubsecdotsep}}
\renewcommand{\cftsecafterpnum}{\vskip2pt}
\setlength{\cftbeforesecskip}{7pt}
\setlength{\cftsubsecindent}{1.4em}
\setlength{\cftsubsubsecindent}{2.9em}

% ---------- Cabeçalho / rodapé ----------
\usepackage{fancyhdr}
\pagestyle{fancy}
\fancyhf{}
\renewcommand{\headrulewidth}{0pt}
\renewcommand{\footrulewidth}{0pt}
\fancyfoot[C]{\footnotesize\color{muted}%
  Realidade de Empresas Produtoras de Soluções de IA\quad\textcolor{line}{\textbar}\quad
  \textcolor{accent}{\thepage}}

% ---------- Links ----------
\usepackage{hyperref}
\hypersetup{
  colorlinks=true, linkcolor=accent, urlcolor=accent, citecolor=accent,
  linktoc=all, breaklinks=true,
  pdftitle={Roteiro de Estudos — Realidade de Empresas Produtoras de Soluções de IA},
  pdfauthor={Tony Esaú de Oliveira}
}

% ---------- columnsep ----------
\setlength{\columnsep}{20pt}
\setlength{\columnseprule}{0.4pt}
\renewcommand{\columnseprulecolor}{\color{line}}
\raggedcolumns % preenche coluna esquerda antes da direita (evita vãos)

% ================= MACROS DE ESTRUTURA =================

% Divisor de Parte (largura total, estilo slide)
\newcommand{\rpart}[3]{%
  \clearpage
  \par\phantomsection\addcontentsline{toc}{section}{#1\hspace{0.5em}#2}%
  \vspace{2pt}%
  \noindent\begin{tcolorbox}[enhanced,colback=panel,colframe=accentDeep,
    boxrule=0pt,leftrule=4pt,arc=2pt,left=14pt,right=14pt,top=12pt,bottom=12pt,
    borderline west={4pt}{0pt}{accent}]
    {\displayfont\fontsize{34}{34}\selectfont\bfseries\color{accent}#1}\;%
    {\displayfont\LARGE\bfseries\color{textcol}#2}\par
    \vspace{4pt}{\color{muted}\itshape #3}
  \end{tcolorbox}\par\vspace{6pt}%
}

% Subseção (1.1) — dentro das colunas
\newcommand{\rsub}[2]{%
  \par\needspace{2.5\baselineskip}\vspace{7pt}%
  \phantomsection\addcontentsline{toc}{subsection}{#1\hspace{0.6em}#2}%
  {\displayfont\large\bfseries\color{accent}#1\hspace{0.5em}#2}\par
  \vspace{1pt}{\color{line}\rule{\linewidth}{0.6pt}}\par\vspace{4pt}%
}

% Subsubseção (1.1.1 / a) — dentro das colunas
\newcommand{\rssub}[2]{%
  \par\needspace{1.8\baselineskip}\vspace{4pt}%
  \phantomsection\addcontentsline{toc}{subsubsection}{#1\hspace{0.5em}#2}%
  {\bfseries\color{accentDim}#1\hspace{0.4em}\color{textcol}#2}\par\vspace{2pt}%
}

% Rótulo de bloco simples (sem TOC), ex.: "Finanças", "Personas de Compra"
\newcommand{\rlabel}[1]{\par\needspace{2.5\baselineskip}\vspace{3pt}%
  {\bfseries\color{accent}#1}\par\vspace{1pt}}

% Linha de apoio (lead) abaixo de um título
\newcommand{\rlead}[1]{{\color{muted}\itshape #1}\par\vspace{3pt}}

% Linha de "Importância" nos blocos de responsabilidades (texto fluido)
\newcommand{\impline}[1]{\par{\bfseries\color{accentDim}Importância:\ }{\color{muted}#1}\par\vspace{5pt}}

% Fonte de tabela / nota pequena
\newcommand{\rfonte}[1]{\par\vspace{2pt}{\footnotesize\color{muted}\itshape #1}\par}

% Imagem com legenda (largura de coluna por padrão) — bloco indivisível
\newcommand{\fig}[4][\linewidth]{%
  \par\medskip\noindent\begin{minipage}{\linewidth}\centering
  {\setlength{\fboxsep}{5pt}\fcolorbox{line}{white}{\includegraphics[width=#1]{#2}}}\par
  \vspace{3pt}%
  {\footnotesize\color{accent}\bfseries #3:\ \normalfont\color{muted}#4\par}%
  \end{minipage}\par\medskip}

% Imagem sem legenda numerada — bloco indivisível
\newcommand{\figplain}[2][\linewidth]{%
  \par\medskip\noindent\begin{minipage}{\linewidth}\centering
  {\setlength{\fboxsep}{5pt}\fcolorbox{line}{white}{\includegraphics[width=#1]{#2}}}%
  \end{minipage}\par\medskip}

% Imagem com legenda descritiva (sem número "Imagem N")
\newcommand{\figcap}[3][\linewidth]{%
  \par\medskip\noindent\begin{minipage}{\linewidth}\centering
  {\setlength{\fboxsep}{5pt}\fcolorbox{line}{white}{\includegraphics[width=#1]{#2}}}\par
  \vspace{3pt}{\footnotesize\color{muted}#3\par}%
  \end{minipage}\par\medskip}

% Caixa de fórmula (destaque verde)
\newcommand{\formula}[2]{%
  \par\medskip\begin{tcolorbox}[enhanced,colback=panelHi,colframe=panelHi,
    boxrule=0pt,arc=3pt,left=12pt,right=12pt,top=8pt,bottom=8pt,
    borderline west={3pt}{0pt}{accent}]
    {\bfseries\color{accent}#1}\\[4pt]{\color{textcol}#2}
  \end{tcolorbox}\par\medskip}

% Início do corpo em duas colunas (atalhos)
\newcommand{\cols}{\begin{multicols}{2}}
\newcommand{\ecols}{\end{multicols}}

% Espaçamento de parágrafo
\setlength{\parindent}{0pt}
\setlength{\parskip}{5pt}

\begin{document}
\color{textcol}
% Quebra de linha tolerante para colunas estreitas sem hifenização PT-BR
\tolerance=2400
\emergencystretch=2.2em
\hyphenpenalty=1200
\hbadness=2400
\sloppy

% =====================================================
% CAPA
% =====================================================
\thispagestyle{empty}
\begin{center}
\vspace*{1.2cm}

\vspace{0.9cm}

{\displayfont\Large\bfseries\color{accent}\addfontfeatures{LetterSpace=8.0}ROTEIRO DE ESTUDOS}

\vspace{0.7cm}

{\color{line}\rule{0.75\linewidth}{0.8pt}}

\vspace{0.7cm}

{\displayfont\fontsize{30}{36}\selectfont\bfseries\color{textcol}
Realidade de Empresas\\ Produtoras de Soluções de\\[12pt] \textcolor{accent}{Inteligência Artificial}}

\vspace{0.7cm}

{\color{line}\rule{0.75\linewidth}{0.8pt}}

\vspace{1.1cm}

\begin{tcolorbox}[enhanced,width=0.82\linewidth,colback=panel,colframe=line,
  boxrule=0.8pt,arc=4pt,left=18pt,right=18pt,top=14pt,bottom=14pt,
  borderline west={3pt}{0pt}{accent}]
  {\bfseries\color{accent}Carga Horária}\hspace{0.6em}\color{textcol}12 horas\par
  \vspace{8pt}
  {\bfseries\color{accent}Conteúdo Programático}\par
  \vspace{4pt}
  \color{textcol}
  \begin{enumerate}[leftmargin=1.7em,itemsep=3pt,topsep=2pt]
    \item Introdução à I.A. aplicada a Empresas (B2B);
    \item Desafios de implementar I.A.;
    \item Como funciona uma empresa produtora de I.A.;
    \item Modelos de negócios;
    \item Dinâmica prática.
  \end{enumerate}
\end{tcolorbox}

\vfill
{\footnotesize\color{muted}Tony Esaú de Oliveira \textcolor{line}{\textbar}\ Residência em I.A. \textcolor{line}{\textbar}\ PPI Softex/SiDi}
\vspace{0.6cm}
\end{center}

% =====================================================
% SUMÁRIO
% =====================================================
\clearpage
\thispagestyle{empty}
{\tableofcontents}
\clearpage

% =====================================================
% PARTE 1
% =====================================================
\rpart{01}{Introdução à I.A. Aplicada a Empresas (B2B)}
{Mercado, ecossistema, oportunidades e como a IA se encaixa na lógica de negócios entre empresas.}

\cols

\rsub{1.1}{O que é B2B?}
\textbf{B2B} (\textit{Business to Business}) é o modelo em que uma empresa vende produtos ou serviços para outra empresa, e não para o consumidor final. Entender esse contexto é a base da disciplina, porque soluções de IA vendidas para empresas seguem essa lógica.

\rssub{a)}{Business to Business}
Empresa vende para empresa.\\
\textbf{Ex.:} fábrica de peças para montadora, software de gestão para outra empresa.

\rssub{b)}{Business to Consumer (B2C)}
Empresa vende direto ao consumidor final.\\
\textbf{Ex.:} supermercado, loja de roupas.

\rssub{c)}{Consumer to Consumer (C2C)}
Consumidor vende para consumidor.\\
\textbf{Ex.:} anúncios entre pessoas físicas em marketplaces.

\rssub{d)}{Business to Government (B2G)}
Empresa vende para o governo, em geral via licitação.

\rssub{1.1.1}{Características do B2B}
O B2B, sigla para \textit{business to business}, é o modelo de negócio em que uma empresa vende produtos ou serviços para outra empresa, e não para o consumidor final, como acontece no B2C. Nesse tipo de relação comercial, quem compra não busca satisfação pessoal, mas sim insumos, ferramentas ou soluções que ajudem a própria empresa a operar, produzir ou revender.

Isso muda bastante a lógica da venda, tornando a decisão mais racional, o relacionamento mais duradouro e a negociação mais técnica. A seguir, estão as principais características que definem esse modelo e o distinguem do comércio voltado ao consumidor.

\fig[0.86\linewidth]{image15.png}{Imagem 1}{Características do B2B (fonte: autor).}

\rsub{1.2}{Mercado de I.A.}
O mercado de inteligência artificial vem crescendo de forma acelerada nos últimos anos e se tornou um dos setores mais estratégicos da economia mundial. Empresas de praticamente todas as áreas, da saúde ao varejo, passando por finanças, indústria e educação, têm investido em soluções de IA para automatizar tarefas, analisar grandes volumes de dados e tomar decisões com mais rapidez e precisão. Essa importância se explica porque a inteligência artificial não é apenas mais uma tecnologia, mas uma ferramenta capaz de transformar a maneira como os negócios funcionam, gerando eficiência, reduzindo custos e criando novos produtos e serviços. Por isso, dominar e acompanhar esse mercado se tornou essencial para quem deseja se manter competitivo no cenário atual.

\begin{itemize}
  \item Agentes de IA devem puxar investimento de \textbf{US\$ 3,4 bilhões} em TI no Brasil em 2026 (IDC, via Mobile Time).
\end{itemize}

\fig[0.78\linewidth]{image18.png}{Imagem 2}{Investimento impulsionado pela I.A. no Brasil (fonte: MEDEIROS, 2026).}

\begin{itemize}
  \item O Brasil está entre os países que mais usam IA, e cerca de \textbf{60\% dos brasileiros} acreditam que ela traz maior expectativa de ganhos (pesquisa Ipsos, via Agência Brasil).
\end{itemize}

\fig[0.78\linewidth]{image17.png}{Imagem 3}{Uso de I.A. (fonte: ALBUQUERQUE, 2025).}

\begin{itemize}
  \item Adoção por país (parcela da população em idade ativa usando IA, H2 2025): \textbf{Emirados Árabes 64\%} e \textbf{Singapura 61\%} lideram; Noruega 46\%, Irlanda 45\%, França 44\%, Espanha 42\%. Os EUA aparecem em \textbf{24º (28\%)} apesar de liderarem em modelos; Brasil em torno de 17\%.
\end{itemize}

\ecols

\figplain[0.98\linewidth]{image5.png}
\begin{center}\vspace{-6pt}{\footnotesize\color{accent}\bfseries Imagem 4:\ }{\footnotesize\color{muted}Adoção de I.A. pelo mundo (fonte: MICROSOFT CORPORATE, 2025).}\end{center}

\cols

\rsub{1.3}{Ecossistema de Inteligência Artificial}
O ecossistema de Inteligência Artificial se organiza em camadas, e cada uma reúne empresas que sustentam a inovação a partir de um ponto diferente da cadeia. Da base, formada pelo hardware, até o topo, onde ficam as aplicações que o usuário final utiliza, todas as camadas se conectam e dependem umas das outras. As principais são:

\begin{itemize}
  \item \textbf{Aplicações de IA generativa:} ficam na ponta mais próxima do usuário e entregam a IA em forma de produto, gerando texto, código, imagens, áudio e vídeo (por exemplo, Synthesia, Grammarly e Runway).
  \item \textbf{Dados:} responsáveis pela coleta, rotulagem, preparação e armazenamento das informações que alimentam os modelos, sem os quais nada funciona (por exemplo, Scale, Databricks, Snorkel e Pinecone).
  \item \textbf{Modelos fundacionais:} os grandes modelos treinados em enormes volumes de dados, que servem de base para diversas aplicações (por exemplo, OpenAI, Google, Anthropic, Meta AI, Mistral e Cohere).
  \item \textbf{Chips:} o hardware especializado que realiza o processamento pesado exigido pela IA, tanto no treino quanto no uso (por exemplo, NVIDIA, AMD, Intel, Qualcomm e Cerebras).
  \item \textbf{Computação e Nuvem:} a infraestrutura que oferece poder de processamento e servidores para treinar e executar os modelos em larga escala (por exemplo, AWS, Google Cloud, Azure, Oracle e IBM).
\end{itemize}

\fig[0.62\linewidth]{image14.png}{Imagem 5}{Ecossistema de I.A. (fonte: VISUAL CAPITALIST, 2025).}

No Brasil, já existe um mapa de mais de 100 startups de IA distribuídas por setores como Saúde, Legal, Financeiro, Agronegócio, Educação, Varejo, RH, Indústria e Governo, além de um amplo guia de ferramentas por finalidade.

\fig[0.72\linewidth]{image6.png}{Imagem 6}{Ranque das Empresas de I.A. no Brasil (fonte: F6S, 2026).}

\rsub{1.4}{Oportunidades}
O gráfico a seguir compara a proporção de tarefas profissionais que os LLMs \textbf{poderiam teoricamente realizar} (área azul) com a \textbf{cobertura observada nos dados de uso reais} (área vermelha). A diferença entre as duas áreas indica o espaço de oportunidade ainda não explorado. Essa diferença tem nome no estudo original da Anthropic: exposição observada. A métrica cruza três fontes, o banco ocupacional americano O*NET, as notas teóricas de viabilidade de Eloundou et al. (2023) e os dados reais de uso do Claude em contexto profissional. O ponto central é o tamanho do vão: em Computação e Matemática, a capacidade teórica chega a 94\%, mas o uso real fica em só 33\%. Esse intervalo é exatamente o espaço de oportunidade ainda não explorado, e a expectativa é que a área vermelha cresça até cobrir a azul conforme a adoção avança.

\fig[0.9\linewidth]{image8.png}{Imagem 7}{Capacidade teórica e observada de agentes de I.A. por categoria (fonte: ANTHROPIC, 2026).}

\rsub{1.5}{Setores-Chave para Adoção}

\rlabel{Finanças}
\begin{itemize}
  \item Automação de processos;
  \item Detecção de fraudes;
  \item Segurança de transações.
\end{itemize}

\rlabel{Saúde}
\begin{itemize}
  \item Melhorar o diagnóstico médico;
  \item Interpretação de exames e imagens;
  \item Suporte e orientação a pacientes.
\end{itemize}

\rlabel{Varejo}
\begin{itemize}
  \item Personalizar a experiência dos clientes;
  \item Recomendações por histórico e preferências;
  \item Análise de compras para ajustar estoque e marketing.
\end{itemize}

\rlabel{Indústria}
\begin{itemize}
  \item Monitorar equipamentos em tempo real;
  \item Evitar falhas e reduzir custos de manutenção;
  \item Aumentar eficiência e automatizar fábricas.
\end{itemize}

\ecols

\rsub{1.6}{Diferenças da IA para B2B e B2C}
\nopagebreak
\begin{center}
\begin{tabularx}{\linewidth}{@{}p{2.3cm} L L@{}}
\rowcolor{accentDeep}\textbf{\textcolor{white}{Critério}} & \textbf{\textcolor{white}{B2C}} & \textbf{\textcolor{white}{B2B}}\\
\textbf{Complexidade} & Negociação de baixa complexidade, ciclo de vendas curto. & Alta complexidade, venda consultiva, ciclo longo (pode durar anos).\\
\rowcolor{panel}\textbf{Volume de Dados} & Muitas transações e interações diárias (Big Data), porém dados menos profundos por usuário. & Menos transações, mas dados mais complexos e estruturados.\\
\textbf{Tomada de Decisão} & Rápida, muitas vezes emocional e individual. & Lenta, racional, envolve múltiplos \textit{stakeholders}.\\
\rowcolor{panel}\textbf{Objetivo} & Aumentar conversões, reduzir abandono de carrinho, fidelização à marca. & Reduzir custos, aumentar produtividade, otimizar o ciclo de vendas longo.\\
\textbf{Exemplos} & Recomendação de produtos (Amazon), algoritmos de conteúdo (TikTok, Netflix). & Análise de chamadas de vendas (Gong.io), CRM preditivo, manutenção preditiva na indústria.\\
\bottomrule
\end{tabularx}
\end{center}

\vspace{10pt}

\cols

\rsub{1.7}{Processo de Compras B2B}

\rlabel{Personas de Compra}
São os diferentes papéis que aparecem numa decisão de compra B2B. Raramente é uma pessoa só, e entender quem é quem ajuda a direcionar a negociação.
\begin{itemize}
  \item \textbf{Decisor:} quem tem autoridade final e aprova a compra, normalmente ligado ao orçamento, como um gestor ou diretor;
  \item \textbf{Influenciador:} quem orienta a escolha com conhecimento técnico ou de mercado, sem necessariamente assinar embaixo, como um especialista ou consultor;
  \item \textbf{Comprador:} quem executa a aquisição na prática, cuidando de cotações, contratos e formalização, em geral o setor de compras;
  \item \textbf{Usuário:} quem vai de fato utilizar a solução no dia a dia e sente na pele se ela resolve o problema, sendo uma voz importante sobre a necessidade real.
\end{itemize}

\rlabel{Termos Comuns}
São as siglas que estruturam as etapas formais de uma compra corporativa, geralmente nesta ordem:
\begin{itemize}
  \item \textbf{RFI (Request for Information):} pedido inicial de informações, usado para conhecer o mercado e levantar quais fornecedores existem e o que oferecem, ainda sem compromisso.
  \item \textbf{RFQ (Request for Quotation):} pedido de cotação, focado em preço e condições, usado quando a empresa já sabe o que quer e compara valores.
  \item \textbf{RFP (Request for Proposal):} pedido de proposta completa, usado em compras mais complexas, em que o fornecedor detalha solução, prazos, metodologia e preço, não só o valor.
  \item \textbf{PO (Purchase Order):} ordem de compra, o documento formal que oficializa a aquisição e funciona como contrato entre as partes, fechando o processo.
\end{itemize}

\ecols

% =====================================================
% PARTE 2
% =====================================================
\rpart{02}{Desafios em Implementar I.A. nas Empresas}
{Por que a adoção nem sempre gera resultado, os principais obstáculos e as estratégias para superá-los.}

\cols

\rsub{2.1}{Percepção dos CEO's}
\begin{itemize}
  \item Apenas \textbf{44\%} dos CEO's são considerados pelos seus CEO's como especialistas em IA;
  \item \textbf{77\%} acreditam que a IA inaugura uma nova era de negócios;
  \item Muitos acham que os principais especialistas de tecnologia da empresa não têm conhecimento e recursos suficientes para apoiar e acelerar resultados.
\end{itemize}

A tensão do slide: enorme urgência (queremos IA agora mesmo) convivendo com falta de clareza sobre o quê, para quê e como.

\fig[0.68\linewidth]{image13.png}{Imagem 8}{Cartum que critica a empolgação sem objetivos sobre a I.A. (fonte: Notícias Tech (Gartner), 2025).}

\rsub{2.2}{Diferentes Realidades}
Empresas de um mesmo setor podem estar em estágios de maturidade muito distintos, de processos manuais e artesanais até linhas totalmente automatizadas. O ponto de partida para adotar I.A. não é o mesmo para todas.

\rsub{2.3}{Resultados Não Percebidos}
Manchetes recentes convergem para um mesmo diagnóstico: a adoção de I.A. cresce, mas boa parte das empresas ainda não gera impacto real.
\begin{itemize}
  \item A adoção dispara, mas a maioria não gera impacto real (McKinsey).
  \item Relatos de que cerca de \textbf{95\% das empresas falham} ao adotar a I.A. e de que \textbf{80\% erram} na implementação.
  \item O Brasil lidera a adoção na América Latina, mas a baixa produtividade ainda é um obstáculo (Fórum Econômico Mundial, potencial de US\$ 1,7 trilhão para a região).
  \item O valor só surge quando \textbf{dados, pessoas e governança} entram no centro da operação.
\end{itemize}
Apesar do entusiasmo e dos investimentos crescentes, a adoção da inteligência artificial nas empresas ainda não se traduz em impacto real na maioria dos casos. Diversos veículos e estudos recentes apontam o mesmo descompasso: a tecnologia se espalha rápido, mas o retorno concreto, em produtividade e resultado, continua modesto, esbarrando em falhas de implementação, ausência de dados organizados, governança e preparo das pessoas.

\ecols

% Imagens lado a lado (2.2 / 2.3)
\begin{center}
\noindent\begin{minipage}{0.48\linewidth}\centering
  {\setlength{\fboxsep}{5pt}\fcolorbox{line}{white}{\includegraphics[width=\linewidth]{image9.png}}}
\end{minipage}\hfill
\begin{minipage}{0.48\linewidth}\centering
  {\setlength{\fboxsep}{5pt}\fcolorbox{line}{white}{\includegraphics[width=\linewidth]{image20.png}}}
\end{minipage}\par
\vspace{3pt}{\footnotesize\color{muted}Diferentes estágios de maturidade em empresas de um mesmo setor.}
\end{center}

\figplain[0.82\linewidth]{image10.png}

\cols

\rsub{2.4}{Desafios que Impedem o ''Agora Mesmo''}
\rfonte{Fonte: Lakshmi Chandrakanth Kasireddy, 2025, p.230.}

\rlabel{Tecnológicos}
\begin{itemize}
  \item Sistemas legados e proprietários;
  \item Processamento e armazenamento de dados;
  \item Disponibilidade de dados;
  \item Segurança.
\end{itemize}

\rlabel{Força de Trabalho e Habilidades}
\begin{itemize}
  \item Escassez de profissionais qualificados;
  \item Temores pela perda do emprego.
\end{itemize}

\rlabel{Financeira e de Recursos}
\begin{itemize}
  \item Altos custos de implantação (HW, SW, RH), com recorrência;
  \item Incerteza de ROI (Retorno sobre o Investimento).
\end{itemize}

\rlabel{Governança, Ética e Regulação}
\begin{itemize}
  \item Viés dos modelos de I.A.;
  \item Falta de transparência;
  \item Potencial uso indevido de I.A.
\end{itemize}

\rsub{2.5}{Estratégias para Superar a Adoção de I.A.}

\rlabel{Deploy de IA faseado}
\begin{itemize}
  \item Minimizar riscos e refinar a estratégia;
  \item Iniciar com pilotos.
\end{itemize}

\rlabel{Capacitar a força de trabalho}
\begin{itemize}
  \item Investir em programas de capacitação;
  \item Esforços colaborativos com a academia.
\end{itemize}

\rlabel{Frameworks de Governança e Ética de I.A.}
\begin{itemize}
  \item Garantir uso responsável;
  \item Transparência, justiça e responsabilidade;
  \item Criar comitês de ética em I.A.
\end{itemize}

\rlabel{Otimizar Custos e ROI}
\begin{itemize}
  \item Soluções em nuvem para escala e eficiência de custos;
  \item Abordagem centrada em dados para maximizar o ROI.
\end{itemize}

\figplain[0.5\linewidth]{image16.png}

\ecols

\rsub{2.6}{Modelo de Maturidade de I.A.}
Um modelo estruturado de maturidade ajuda a organização a identificar em que ponto está e quais etapas levam ao próximo estágio. A tabela a seguir sintetiza os estágios da maturação da implementação de I.A., suas características e as ações recomendadas a cada fase.

\begin{center}
\begin{tabularx}{\linewidth}{@{}>{\bfseries\RaggedRight\arraybackslash}p{3.1cm} L L@{}}
\rowcolor{accentDeep}\textcolor{white}{Estágio} & \textbf{\textcolor{white}{Características}} & \textbf{\textcolor{white}{Ações Recomendadas}}\\
Experimental & Adoção limitada a projetos piloto. & Identificar casos de uso de alto valor, criando POC's (Provas de Conceito) de I.A.\\
\rowcolor{panel}Operacional & I.A. incorporada em alguns fluxos de trabalho. & Padronizar os fluxos de I.A. e treinar a força de trabalho.\\
Transformacional & I.A. integrada em várias funções de negócio. & Automatizar processos e estabelecer governança de I.A.\\
\rowcolor{panel}Toda a empresa & I.A. é fundamental para a estratégia e a tomada de decisões. & Dimensionar a infraestrutura e otimizar modelos de negócio movidos por I.A.\\
\bottomrule
\end{tabularx}
\end{center}
\rfonte{Fonte: Lakshmi Chandrakanth Kasireddy, 2025, p.235.}

\cols

\rsub{2.7}{Modelo de Implementação de I.A. (8 passos)}
\begin{enumerate}
  \item Definir metas;
  \item Avaliar dados;
  \item Escolher tecnologia;
  \item Criar equipe proficiente;
  \item Promover cultura de IA;
  \item Gerenciar riscos;
  \item Testar e avaliar modelos;
  \item Planejar escala e melhoria contínua.
\end{enumerate}
\rfonte{Fonte: Stryker (IBM Think), 2025.}

\figplain[0.82\linewidth]{image2.png}

\rsub{2.8}{Implementar KPI's (Indicador-Chave de Desempenho)}
Medir I.A. exige aplicar a lógica clássica de ROI a um contexto mais complexo. São 4 dimensões e 16 indicadores:
\begin{enumerate}
  \item \textbf{Investimento (tamanho e eficiência)}
    \begin{itemize}
      \item Custo total de I.A.;
      \item Custo por inferência;
      \item Custo por usuário ativo;
      \item Investimento em infraestrutura.
    \end{itemize}
  \item \textbf{Operacional (impacto no trabalho)}
    \begin{itemize}
      \item Horas poupadas;
      \item Tarefas automatizadas;
      \item Redução de retrabalho;
      \item Tempo de ciclo reduzido.
    \end{itemize}
  \item \textbf{Financeira (tradução para o negócio)}
    \begin{itemize}
      \item ROI;
      \item Receita atribuída à I.A.;
      \item Payback;
      \item Eficiência econômica marginal.
    \end{itemize}
  \item \textbf{Desempenho Técnico}
    \begin{itemize}
      \item Acurácia e precisão;
      \item Taxa de automação bem-sucedida;
      \item Latência média;
      \item Adoção e engajamento.
    \end{itemize}
\end{enumerate}

\rsub{2.9}{Casos de Sucesso}

\rssub{2.9.1}{Agronegócio · Cartografia e Mapeamento com Dados Públicos}
\rlead{Ciência de dados e visão computacional para roteamento de máquinas autônomas.}
\textbf{Desafios:} ineficiência no roteamento das máquinas, baixa confiabilidade na atualização de ruas e rotas por GPS, desperdício de implementos agrícolas.

\textbf{Resultados:} automatização da geração de coordenadas, algoritmo de orientação e localização mais precisos, download automático de coordenadas, redução de falhas e roteamentos redundantes, processamento de imagens de satélite para identificar limites de ruas e talhões.

\figplain[\linewidth]{image19.png}

\rssub{2.9.2}{Energia · Sistema Especialista de Operação de Comporta de Barragens}
\rlead{Ciência de dados, NLP e visão computacional para análise de riscos e apoio à decisão.}
\textbf{Desafios:} prevenir impactos de desastres naturais, reduzir a subjetividade na decisão, extrair o conhecimento tácito, consolidar informações de fontes diversas, otimização energética.

\textbf{Resultados:} assistente virtual de apoio à decisão, rastreabilidade e evolução no controle de cheias, maior segurança operacional e para o operador, mitigação do impacto socioambiental.

\figcap[\linewidth]{image11.png}{Sistema especialista de operação de barragens (fonte: SiDi).}

\ecols

% =====================================================
% PARTE 3
% =====================================================
\rpart{03}{Como Funciona uma Empresa Produtora de I.A.}
{Papéis do time, responsabilidades e processos de desenvolvimento.}

\rsub{3.1}{O Time Típico}
Uma empresa que produz soluções de IA reúne frentes técnicas, de produto e comerciais, além de uma frente transversal de segurança e jurídico. Todas jogam juntas para levar a I.A. do laboratório até o cliente.

\figplain[0.86\linewidth]{image7.png}

\begin{center}
\begin{tabularx}{\linewidth}{@{}>{\bfseries\RaggedRight\arraybackslash}p{3.4cm} L@{}}
\rowcolor{accentDeep}\textcolor{white}{Frente} & \textbf{\textcolor{white}{Papéis do Time}}\\
Técnico e Engenharia & \normalfont Engenharia de IA · Engenharia de Dados · MLOps · Engenharia de Software / Plataforma · Arquiteto de Software.\\
\rowcolor{panel}Produto e Experiência & \normalfont Gestão de Produto · UX / Design de Produto.\\
Comercial e Mercado & \normalfont Vendas · Customer Success · Marketing.\\
\rowcolor{panel}Transversal & \normalfont Segurança / Jurídico.\\
\bottomrule
\end{tabularx}
\end{center}

\rsub{3.2}{Responsabilidades}
\cols
Numa empresa de inteligência artificial, entregar valor não depende de uma única função, mas da coordenação entre diferentes áreas que atuam em etapas complementares. O trabalho costuma se dividir entre quem cuida dos dados, garantindo coleta, limpeza e organização, quem desenvolve e treina os modelos, quem cuida da infraestrutura e da colocação desses modelos em produção, e quem conecta tudo isso às necessidades do negócio e do cliente. A essas frentes se somam papéis de governança, ética e segurança, cada vez mais centrais para assegurar que as soluções sejam confiáveis, transparentes e adequadas às regras do mercado. Compreender essa divisão de responsabilidades ajuda a enxergar por que o sucesso de um projeto de I.A. raramente é fruto de uma pessoa só, e sim de uma cadeia bem articulada de competências técnicas e estratégicas.

\rssub{3.2.1}{Responsabilidades de Negócio}

\rlabel{Vendas}
\begin{itemize}\item Elaborar propostas comerciais, RFPs e apresentações executivas.\item Alinhar expectativas técnicas e de negócio.\end{itemize}
\impline{Sustentabilidade financeira da empresa.}

\rlabel{Customer Success}
\begin{itemize}\item Acompanhar \textit{onboarding}, implementação e adoção.\item Monitorar o uso da solução e sugerir melhorias.\item Suportar integrações e resolver incidentes.\item Ajudar a transformar pilotos em contratos escaláveis.\end{itemize}
\impline{Garantir retenção, \textit{upsell} e longevidade dos contratos.}

\rlabel{Arquiteto de Software}
\begin{itemize}\item Apoiar a área de Vendas.\item Traduzir as necessidades do cliente em arquitetura de software.\end{itemize}
\impline{Orquestrar vendas, cliente e times técnicos.}

\rlabel{Gestão de Produto}
\begin{itemize}\item Identificar oportunidades de I.A. nos processos dos clientes.\item Traduzir requisitos de negócio em requisitos técnicos.\item Garantir que o produto entregue valor e reduza dores corporativas.\end{itemize}
\impline{Evita que a empresa construa ``I.A. pela I.A.'', direcionando o esforço para o que realmente gera valor empresarial.}

\rlabel{Marketing}
\begin{itemize}\item Produzir estudos, \textit{whitepapers} e casos de uso.\item Gerar demanda (\textit{leads}) com foco em decisores corporativos.\item Educar o mercado e diferenciar a solução da concorrência.\end{itemize}
\impline{Ajudar as empresas a entender o que a I.A. faz de verdade é parte fundamental da venda.}

\rssub{3.2.2}{Responsabilidades Técnicas}
As frentes técnicas sustentam o núcleo do produto: modelos, dados e operação em produção.

\rlabel{Engenharia de I.A.}
\begin{itemize}\item Desenvolver modelos supervisionados, não supervisionados ou generativos.\item Fazer \textit{fine-tuning} de modelos pré-treinados (LLMs, modelos de visão etc.).\item Criar pipelines de treinamento e validação.\item Otimizar desempenho, latência e custo das inferências.\item Traduzir requisitos de negócio em soluções técnicas.\end{itemize}
\impline{Área que entrega o núcleo tecnológico que diferencia a empresa no mercado.}

\rlabel{Engenharia de Dados}
\begin{itemize}\item Construir \textit{pipelines} ETL/ELT para ingestão, limpeza e transformação de dados.\item Criar e manter \textit{data lakes} e \textit{data warehouses}.\item Padronizar formatos e metadados para uso pelos cientistas de ML.\item Integrar dados de ERPs, CRMs, sensores IoT e bases externas.\item Garantir qualidade (QA), \textit{discoverability} e governança dos dados.\end{itemize}
\impline{Garante dados acessíveis, confiáveis e seguros. Sem dados bem preparados nenhum modelo funciona, e esse é um dos maiores gargalos nas empresas.}

\rlabel{MLOps}
\begin{itemize}\item Orquestrar pipelines de CI/CD para modelos de IA.\item Monitorar modelos em produção (drift, performance, custos).\item Gerenciar conteinerização (Docker), servidores de inferência e GPUs.\item Estabelecer versionamento de dados e modelos.\item Automatizar re-treinamentos periódicos e validações contínuas.\end{itemize}
\impline{É o elo entre pesquisa e operação. Sem MLOps, o modelo funciona no Jupyter Notebook, mas não escala para clientes reais.}

\rssub{3.2.3}{Responsabilidades de Produto}

\rlabel{UX / Design de Produto}
\begin{itemize}\item Projetar \textit{dashboards}, \textit{workflows} e interfaces intuitivas.\item Criar experiências apropriadas para automação com I.A. (chatbots, copilots, assistentes).\item Garantir que a I.A. explique suas recomendações (\textit{explainability}).\item Fazer testes com usuários (\textit{usability tests}).\end{itemize}
\impline{A adoção pelo cliente depende muito da facilidade de uso.}

\rlabel{Engenharia de Software / Plataforma}
\begin{itemize}\item Criar API's que expõem os modelos para os clientes.\item Integrar a I.A. com sistemas legados.\item Criar dashboards, painéis e plataformas web.\item Garantir escalabilidade, resiliência e uso eficiente de recursos.\end{itemize}
\impline{Constrói a aplicação que operacionaliza a I.A. A I.A. sozinha não vira produto, ela precisa de \textit{software} para chegar ao cliente.}

\rlabel{Segurança / Jurídico}
\begin{itemize}\item Avaliar riscos de I.A. (viés, drift, segurança, privacidade).\item Garantir adequação a normas como LGPD, A.I. Act e frameworks do NIST.\item Avaliar contratos e termos de uso.\item Realizar auditorias dos pipelines de I.A.\end{itemize}
\impline{Empresas só compram I.A. se ela for segura, auditável e confiável.}
\ecols

\rsub{3.3}{Processos para Desenvolvimento de Soluções}
Há dois caminhos típicos. No modelo de produto, a empresa itera um MVP até encontrar valor e escala. No modelo consultivo, o trabalho segue um funil comercial, do primeiro contato até a venda.

\phantomsection\addcontentsline{toc}{subsubsection}{3.3.1\hspace{0.5em}Produto (MVP): Ciclo de Venture Design}%
\phantomsection\addcontentsline{toc}{subsubsection}{3.3.2\hspace{0.5em}Projetos (Consultivo): Funil Comercial}%
\begin{center}
\noindent\begin{minipage}[t]{0.58\linewidth}\centering
  {\bfseries\color{accentDim}3.3.1\ \color{textcol}Produto (MVP): Ciclo de Venture Design}\par\vspace{5pt}
  {\setlength{\fboxsep}{5pt}\fcolorbox{line}{white}{\includegraphics[width=\linewidth]{image12.png}}}\par
  \vspace{3pt}{\footnotesize\color{accent}\bfseries Imagem 9:\ \normalfont\color{muted}Processo de Venture Design (fonte: COWAN, 2026).}
\end{minipage}\hfill
\begin{minipage}[t]{0.38\linewidth}\centering
  {\bfseries\color{accentDim}3.3.2\ \color{textcol}Projetos (Consultivo): Funil Comercial}\par\vspace{5pt}
  {\setlength{\fboxsep}{5pt}\fcolorbox{line}{white}{\includegraphics[width=\linewidth]{image1.png}}}\par
  \vspace{3pt}{\footnotesize\color{accent}\bfseries Imagem 10:\ \normalfont\color{muted}Funil Comercial (fonte: autor).}
\end{minipage}
\end{center}

% =====================================================
% PARTE 4
% =====================================================
\rpart{04}{Modelos de Negócios para Projetos de I.A.}
{Como a empresa cobra, constrói soluções sob medida e monta propostas.}

\rsub{4.1}{Modelo de Negócios}
\begin{center}
\begin{tabularx}{\linewidth}{@{}>{\bfseries\RaggedRight\arraybackslash}p{4cm} L@{}}
\rowcolor{accentDeep}\textcolor{white}{Modelo} & \textbf{\textcolor{white}{Descrição}}\\
IA como serviço (AIaaS) & \normalfont Permite usar I.A. sem investir pesado na tecnologia subjacente. Os recursos são oferecidos por uma plataforma em nuvem, via assinatura. Gera receita recorrente ao provedor e dá flexibilidade a quem quer usar I.A. sem compromisso de longo prazo.\\
\rowcolor{panel}Preços por Resultado e Criação de Valor & \normalfont O pagamento fica ligado ao desempenho e aos resultados da solução, não às funcionalidades. Alinha os incentivos entre fornecedor e cliente e estimula a melhoria contínua e o ajuste fino das aplicações de I.A.\\
Soluções Personalizadas & \normalfont Base dos modelos de consultoria. A empresa cria sistemas de I.A. sob medida, alinhados às necessidades operacionais e aos objetivos estratégicos de cada cliente.\\
\bottomrule
\end{tabularx}
\end{center}

\cols
\rlabel{Exemplos de Mercado}
\begin{itemize}
  \item e-core Orbit AI (e-core.com/pt-br/orbit-ai);
  \item Google Cloud, para modernização de infraestrutura em nuvem;
  \item IBM Watson;
  \item Fin, agente de I.A. para atendimento ao cliente.
\end{itemize}

\rsub{4.2}{Construção de Soluções Customizadas}
As três primeiras fases descobrem o problema certo e checam se vale a pena resolvê-lo.

\rlabel{Necessidades do Cliente}
\begin{itemize}\item Levantar o problema.\item Identificar tecnologias e habilidades necessárias.\end{itemize}

\rlabel{Definir o Problema}
\begin{itemize}\item Qual oportunidade será resolvida.\item Definir quem é atendido pela solução.\end{itemize}

\rlabel{Avaliar Viabilidade}
\begin{itemize}\item Considerar recursos disponíveis, capacidades técnicas e alinhamento estratégico.\item Equilibrar desejo, viabilidade e possibilidades.\end{itemize}

\textbf{Ferramentas:} entrevistas com \textit{stakeholders}, jornada da experiência do usuário, Matriz CSD, Jobs to be done e Matriz SWOT.

\figplain[\linewidth]{image3.png}

As três fases seguintes traduzem a solução em valor, diferencial e retorno financeiro:

\rlabel{Modelar Valor}
\begin{itemize}\item Mapear como a oportunidade gera valor para o cliente e para a empresa.\item Identificar benefícios e impactos financeiros.\end{itemize}

\rlabel{Diferencial Competitivo}
\begin{itemize}\item Definir a proposta de valor e o diferencial (\textit{unfair advantage}).\item Entender como as tendências impactam solução e negócio.\end{itemize}

\rlabel{Determinar o ROI}
\begin{itemize}\item Calcular o ROI potencial (custos de desenvolvimento, implementação e manutenção).\item Comparar com outras oportunidades para priorizar.\end{itemize}

\textbf{Ferramentas:} Business Model Canvas, análise de benchmark e Matriz SWOT.

\formula{Cálculo de ROI}{$\displaystyle \mathrm{ROI}\,(\%) \;=\; \frac{\text{receita} - \text{custo}}{\text{custo}} \times 100$}

\rsub{4.3}{Proposta Técnico-Comercial}
Uma proposta técnico-comercial é o documento que uma empresa apresenta a um cliente para formalizar uma oferta, reunindo, num só lugar, tanto a solução que será entregue quanto às condições comerciais dessa entrega. Do lado técnico, ela descreve o escopo do trabalho, a metodologia, as tecnologias envolvidas, os prazos e os resultados esperados, mostrando que a empresa entende o problema e sabe como resolvê-lo. Do lado comercial, ela detalha preços, formas de pagamento, garantias e demais termos do acordo.

\rlabel{Características}
\begin{itemize}\item Centrada no cliente, alinhada aos seus desafios.\item Traz informações válidas para diferentes decisores.\item É objetiva, clara e sucinta.\item Tem visual profissional.\item É entregue sem erros de gramática e de cálculo.\end{itemize}

\rlabel{Estrutura}
\begin{itemize}\item Capa.\item Visão geral.\item Objetivos.\item Escopo e não escopo.\item Entregáveis.\item Cronograma.\item Por que escolher você.\item Investimento e condições de pagamento.\end{itemize}

\rlabel{Exemplos de Proposta}
\begin{itemize}
  \item Modelo de proposta comercial para aplicativo móvel;
  \item Modelo de proposta de desenvolvimento de software (Prospero);
  \item Proposta paipe-software.
\end{itemize}

\rssub{4.3.1}{Cálculo de Esforço Técnico}
O cálculo de esforço prevê o tempo (horas por pessoa) e os recursos necessários para concluir um projeto. Uma divisão comum é 60\% em desenvolvimento, 20\% em testes e de 10\% a 20\% em gestão.

\rlabel{Passos para Calcular}
\begin{itemize}
  \item Definir o escopo, detalhando funcionalidades e requisitos;
  \item Quebrar em tarefas (WBS), dividindo o projeto em atividades menores e gerenciáveis;
  \item Estimar horas por tarefa, aplicando técnicas de estimativa;
  \item Considerar fatores de produtividade, como pausas, reuniões e interrupções (a eficiência real é menor que 8 horas por dia);
  \item Adicionar margem de risco de 10\% a 20\% para imprevistos.
\end{itemize}

\figplain[\linewidth]{image4.png}

\rlabel{Técnicas de Estimativa}
\textbf{\textcolor{accent}{Paramétrica.}} Baseia-se em dados históricos. Exemplo: um dia para desenvolver 1.000 linhas de código.\par\smallskip
\textbf{\textcolor{accent}{PERT.}} Calcula uma média ponderada de três cenários (otimista, mais provável e pessimista).\par\smallskip
\textbf{\textcolor{accent}{Pontos de Função.}} Mede a funcionalidade entregue ao usuário e multiplica pela produtividade da equipe (ex.: 500 PF × 10 horas por PF).\par\smallskip
\textbf{\textcolor{accent}{Planning Poker (ágil).}} A equipe atribui pontos de complexidade (Fibonacci) às tarefas e depois os traduz em esforço.

\formula{Estimativa PERT}{$\displaystyle X \;=\; \frac{P + 4\cdot MP + O}{6}$\\[5pt]{\small\color{muted}(P pessimista, MP mais provável, O otimista)}}
\ecols

% =====================================================
% PARTE 5
% =====================================================
\rpart{05}{Dinâmica Prática}
{Hands-On: construir, do zero ao repositório, uma proposta técnico-comercial completa.}

\cols
Construir uma proposta técnico-comercial completa para um cliente fictício, aplicando o processo de solução customizada, o cálculo de esforço e a estrutura de proposta, e deixar tudo versionado no repositório como artefato de portfólio.

\begin{enumerate}
  \item \textbf{Escolher o Cenário:} definir o problema de negócio que resolve (por exemplo, previsão de demanda no varejo, automação de atendimento, visão computacional para controle de qualidade). Escrever em uma ou duas frases quem é o cliente fictício e qual a dor dele;
  \item \textbf{Entender o Problema e Definir Escopo:} registrar os requisitos, as restrições e as premissas. Deixar explícito o que a solução vai contemplar e, igualmente importante, \textit{o que não vai}. Esse recorte de escopo é o que sustenta o cálculo de esforço depois;
  \item \textbf{Desenhar a Solução Customizada:} descrever a abordagem técnica, as tecnologias envolvidas, as etapas do projeto e a arquitetura geral da solução, mostrando como ela resolve a dor definida no passo 1;
  \item \textbf{Calcular o Esforço:} estimar as etapas, os perfis envolvidos (dados, modelagem, infraestrutura, gestão), o tempo de cada atividade e o custo correspondente. Guardar isso preferencialmente numa planilha, com as premissas de cada estimativa anotadas, para dar rastreabilidade ao número final;
  \item \textbf{Montar a Proposta:} organizar o documento nas seções esperadas, ou seja, contexto e entendimento do problema, solução proposta, escopo, metodologia e cronograma, equipe, investimento e condições comerciais, e diferenciais;
  \item \textbf{Publicar em Repositório:} subir a proposta (em Markdown ou PDF), um README explicando o cenário e as premissas, e uma pasta de anexos com a planilha de esforço. Ir registrando a evolução em commits, para o repositório mostrar o processo e não só o resultado final.
\end{enumerate}

\rlabel{Estrutura de pastas sugerida para o repositório}
\begin{itemize}
  \item \texttt{README.md}, com a visão geral do caso e as premissas;
  \item \texttt{/proposta}, com o documento da proposta técnico-comercial;
  \item \texttt{/esforco}, com a planilha de cálculo de esforço;
  \item \texttt{/anexos}, com material de apoio (diagramas, referências).
\end{itemize}
\ecols

% =====================================================
% REFERÊNCIAS
% =====================================================
\clearpage
\vspace{6pt}
\phantomsection\addcontentsline{toc}{section}{Referências}
\noindent{\displayfont\LARGE\bfseries\color{accent}Referências}\par
\vspace{2pt}{\color{line}\rule{\linewidth}{0.8pt}}\par\vspace{6pt}

{\footnotesize\setlength{\parindent}{0pt}%
\newcommand{\rref}[1]{\par\hangindent=1.4em\hangafter=1 #1\par\vspace{4pt}}
\cols
\rref{AIBRASIL. \textbf{``Guia de Ferramentas de IA''}. AI Brasil (LinkedIn), 2026. Acesso em: 12/03/2026.}
\rref{ALBUQUERQUE, Flávia. \textbf{``Brasil está entre os países que mais usam inteligência artificial''}. Agência Brasil, 2025. Acesso em: 26/02/2026.}
\rref{ANTHROPIC, C. \textbf{``Labor market impacts of AI: A new measure and early evidence''}. Anthropic, 2026. Acesso em: 11/03/2026.}
\rref{CHUNG, Wan Ming. \textbf{``As 4 dimensões e os 16 indicadores que mostram se a sua IA realmente entrega valor''}. Pulsar WCA Newsletter (LinkedIn), 13/02/2026. Acesso em: 26/02/2026.}
\rref{F6S. \textbf{``100 top AI (Artificial Intelligence) companies and startups in Brazil in February 2026''}. F6S, 2026. Acesso em: 26/02/2026.}
\rref{LAKSHMI CHANDRAKANTH KASIREDDY. \textbf{``Overcoming Adoption Barriers: Strategies for Scalable AI Transformation in Enterprises''}. Journal of Informatics Education and Research, v. 5, n. 2, 2025. DOI: 10.52783/jier.v5i2.2459. Acesso em: 26/02/2026.}
\rref{MEDEIROS, Henrique. \textbf{``Agentes de IA vão puxar investimento de US\$ 3,4 bilhões em TI no Brasil em 2026''}. Mobile Time, 2026. Acesso em: 26/02/2026.}
\rref{NEUFELD, Dorothy. \textbf{``Mapped: AI Adoption Rates by Country''}. Visual Capitalist, 2025. Acesso em: 26/02/2026.}
\rref{NOTÍCIAS TECH. \textbf{``Gartner: CEOs acreditam que suas equipes executivas não têm conhecimento elevado em IA''}. Notícias Tech, 2025. Acesso em: 26/02/2026.}
\rref{STRYKER, Cole. \textbf{``Implementação de inteligência artificial: 8 passos para o sucesso''}. IBM Think, 2025. Acesso em: 26/02/2026.}
\rref{VISUAL CAPITALIST. \textbf{``Visualized: The AI Ecosystem''}. Visual Capitalist, 2025. Acesso em: 26/02/2026.}
\rref{COWAN, Alex. \textbf{Venture Design Process}. Cowan+, 2026. Disponível em: alexandercowan.com/venture-design.}
\ecols
}

\end{document}
